% =======================================
% MEETING 10
% ---------------------------------------
% DATE   - September 23, 2026
% FORMAT - F2F
% =======================================
% TOPIC  - Metric Tensor
% =======================================
\subsection{The Metric Tensor}

	The metric tensor is the most important algebraic object in General relativity.
	It allows is to intrinsically define the curvature, lengths, and angles, inside a 
	general manifold. 

\subsubsection{Definition}
	We define the metric as such
	\begin{definition}{Metric Tensor}
		A metric tensor, $g\in V^*\times V^*$, is a symmetric, nondegenerate $(0,2)$ tensor acting on 
		$V$.
	\end{definition}
	

	First, some qualifiers. Symmetric means, as we have previously discussed, is that, for all vectors 
	$u,v\in V$,\footnote{this is now a tensor acting on two vectors, despite the $(\cdot)$}
	\[
		g(u,v) = g(v,u)
	\]
	Which implies
	\[
		g_{\mu\nu} = g_{\nu\mu}
	\]
	which is also equivalent to
	\[
		 g(\p_\mu,\p_\nu) = g(\p_\nu, \p_\mu)
	\]

	Nondegenerate means that if there exists a vector $v\in V$, such that
	\[
		g(u,v) = 0
	\]
	for all $u$. Then this can only happen if $v = 0$. Otherwise, this is a degenerate tensor.

	Tersely, we have
	\[
		g(v,u) = 0\forall u \in V\implies v=0\in V
	\]

	This is in the same sense as a matrix degeneracy. In the matrix isomorphism
	of $g$, $g_{\mu\nu}$, we require that the metric has a nonzero determinant
	\[
		\det(g_{\mu\nu})\neq 0
	\]
	If that happens at any point, then the coordinates are wrong for that point.
	That's why in spherical coordinates, $r=0$ and the poles $\phi = 0$ and $\phi = \pi$,
	are not valid.

	We have what we call a volume element, $\text{vol}$, which is a 3-form. In the Euclidean metric,
	which is an identity metric (also symmetric)
	\[
		\tensor{e}{_{ij}} = \bmqty{
			1 & 0 & 0 \\
			0 & 1 & 0 \\
			0 & 0 & 1 \\
	} 
	\]
	The volume element is just 
	\[
		\text{vol} = dx^x dx^y dx^z
	\]
	However, for a general metric, the volume element is
	\[
		\text{vol} = \sqrt{\det(g_{\mu\nu}) dx^\mu}
	\]
	The nondegenerate requirement prevents the volume element from vanishing.
	The volume element must not be zero for any vector.

	Second, as we have noted before, the metric tensor is an $(0,2)$ tensor. That means it acts 
	on two vectors in the vector space $V$. 
	
	Naturally, since it is a tensor, it is represented in basis form as 
	\[
		g = g_{\mu\nu} dx^\mu dx^\nu
	\]
	Where $dx^\mu dx^\nu$ is implied to be a symmetric tensor product. To find the metric tensor
	components, we simply feed it the coordinate basis vectors on the manifold $\mcal M$, just like
	how we find the components of a general tensor.
	\[
		g_{\mu\nu}  = g(\p_\mu, \p_\nu)
	\]
	Later on, we will define those components more concretely. 

	The metric tensor allows us to take two vector and relate them to one another via a number.
	We shall see how the metric is useful on the next few chapters.

\subsubsection{Inner Product}
	\label{subsubsec:inner-product}
	We mentioned several lectures ago that
	to get the notion of lengths and orthogonality
	between vectors, we need an operation that will
	allow us to measure lengths and angles. An operation
	which takes two vectors and spits out a scalar.
	We shall discuss that operation now.

	The inner product is an operation that takes two vectors, operates on them,
	and returns a real number. How do we make sense of that number?

	We can define an inner product as the following.
	\begin{definition}{Inner Product}
		For two vectors $u,v\in T\mcal{M}$ with respect to the metric $g$,
		the inner product is defined as
		\begin{equation}
			\label{eq:inner-product}
			u\cdot v = g(u,v).
		\end{equation}

		The inner product is a map such that
		\[
			g:T_p\mcal{M}\times T\mcal{M}\to\R.
		\]
	\end{definition}

	The inner product then describes the \textit{relationship} of one vector
	to another. Thus, it is a natural way to also describe lengths. But what is a length?
	It is most natural to think of the length of a vector as the magnitude of the vector.
	In classical vector physics, the magnitude is the square root of the dot product
	of a vector with itself, since we define that $v^2 = v\cdot v$. There is a
	similar notion in differential geometry.

	A metric is similar to a familiar inner product in matrices, that is, the dot product.
	However, the difference between them is that $g(u,v)$ can be negative, and $g(u,v)=0$
	does not mean $v=0$. Once a matric is defined for a vector space $V$, we can define the 
	lengths and orthogonality between its elements. 

	We then define the following:
	\begin{definition}{Vector Magnitude}
		The length, magnitude, or \textit{norm}, $\norm{v}$ of a vector $v\in V$, given a metric
		$g$, is defined as
		\begin{equation}
			\label{eq:magnitude}
			\norm{v} = \sqrt{\abs{g(v,v)}}
		\end{equation}
	\end{definition}
	We say $\abs{g(v,v)}$ because, for a general metric $g$, the dot product 
	may be negative. This can also be represented as $\sqrt{\pm g(v,v)}$.

	Why does this give us a notion of length? In Euclidean space, the
	Pythagorean theorem tells us that the square of the length of a vector
	is the sum of the squares of its components. The metric generalizes
	this notion to an arbitrary manifold. In a coordinate basis,
	\[
		g(v,v)
		=
		g_{\mu\nu}v^\mu v^\nu.
	\]
	Thus, $g(v,v)$ plays the role of the squared length of the vector,
	while $\sqrt{g(v,v)}$ gives its length.

	What other notions does the usual dot product give us? It allows us to find
	the angles between vectors. If $v\cdot u=0$, then $v$ and $u$ are orthogonal.
	We say that vectors are orthogonal if their inner product is zero. If it 
	is not, they have components that align to each other. If the dot product 
	between them is maximal, then that means the vectors are aligned, or coincidental.
	For unit normal vectors, the maximal dot product is simply $v\cdot v = 1$.
	We say that a set of basis vectors is orthonormal if the basis vectors
	are unit vectors and orthogonal to one another.

	If the basis vectors are not orthonormal, then the inner products
	between them need not be zero or one. If the basis vectors are
	orthonormal, their inner products satisfy
	\[
		e_\mu\cdot e_\nu = \delta_{\mu\nu}.
	\]

	We can define this orthonormality more generally for any 
	manifold, not just flat space.

	\begin{definition}{Orthonormality}
		We say that a set of basis vectors
		$\{e_\mu\}$ is orthonormal if, for all $e^\mu \in\{e_\mu\}$, 
		\begin{equation}
			\label{eq:orthonormality}
			g(e_\mu,e_\nu)=\pm \delta_{\mu\nu}.
		\end{equation}
		That is,
		\[
			g(e_\mu,e_\mu) = \pm 1
		\]
	\end{definition}
	
	Once again, we say $\pm$ because, in general, the inner product may 
	not be positive.

	In the Euclidean manifold, the inner product is simply the familiar dot product.
	Here, we recognize that the dot product, in tensor notation, is just a
	tensor application. Specifically it is called the Euclidean metric.
	\begin{align*}
		e_{ij}
		&=
		e_{\mu\nu}dx^\mu\otimes dx^\nu\\
		&=\delta_{\mu\nu}dx^\mu\otimes dx^\nu
	\end{align*}
	acting on the vectors $\vb{u}$ and $\vb{v}$.
	\begin{align*}
		\vb{u}\cdot\vb{v}
		&= e(u,v)\\
		&= e(u^\mu\p_\mu,v^\mu\p_\nu)\\
		&= u^\mu v^\nu e(\p_\mu,\p_\mu)\\
		&= u^\mu v^\nu \delta_{\mu\nu} \\
		&= u^\mu v_\mu 
	\end{align*}
	Which is just the sum of the product of the components. It's a regular, garden variety
	dot product.

	We can also see this directly from the tensor expression for the metric:
	\[
		e
		=
		\delta_{\mu\nu}dx^\mu\otimes dx^\nu.
	\]
	Since
	\[
		dx^\mu(\p_\nu)=\delta^\mu_\nu,
	\]
	we have
	\begin{align*}
		e(\p_\mu,\p_\nu)
		&=
		\delta_{\alpha\beta}
		dx^\alpha(\p_\mu)
		dx^\beta(\p_\nu)\\
		&=
		\delta_{\alpha\beta}
		\delta^\alpha_\mu
		\delta^\beta_\nu\\
		&=
		\delta_{\mu\nu}.
	\end{align*}

	We also recognize that the dot product in the Euclidean metric provides
	the notion of angles between vectors. In particular,
	\[
		\vb{u}\cdot\vb{v}
		=
		\abs{\vb{v}}\abs{\vb{u}}\cos\theta.
	\]
	Therefore, for a general metric tensor,
	\[
		\cos\theta
		=
		\frac{g(u,v)}
		{\sqrt{g(u,u)}\sqrt{g(v,v)}}.
	\]
	We say then that the vectors are orthogonal, if 
	\[
		g(u,v) = 0
	\]
	\paragraph{Metric Types}
		By the definition of an orthonormal basis, i.e.,
		\[
			g_{\mu\nu} = \tensor{\delta}{_{ij}}
			\begin{cases}
				0, &\mu\neq\nu\\
				\pm 1,&\mu=\nu
			\end{cases}
		\]
		then an orthonormal basis's matrix representation is a diagonal matrix.
		\[
			g_{\mu\nu}
			=
			\bmqty{
				\pm 1      & 0          & \cdots & 0 \\
				0          & \pm 1      & \cdots & 0 \\
				\vdots     & \vdots     & \ddots & \vdots \\
				0          & 0          & \cdots & \pm 1
			}
		\]

		Beyond having a diagonal matrix in an orthonormal basis, metrics 
		whose diagonal elements are all $+1$ are said to be \textbf{positive definite}
		or \textbf{Riemannian}. An example is the Euclidean metric. 

		Metrics whose 
		diagonal elements are all $-1$ are said to be \textbf{negative definite},
		and others are said to be \textbf{indefinite} or \textbf{semi-Riemannian}. 
		The indefinite metrics whose diagonal elements only have one $-1$, i.e., $g(\p_\mu,\p_\mu)=-1$
		for one specific $\mu$, then it is called a \textbf{Lorentzian} metric. 

		The one most used in relativity are Lorentzian metrics and positive definite metrics.
		\textbf{The existence of a Lorentzian metric on a manifold defines a spacetime}. The very definition
		of a spacetime (Def. \ref{def:spacetime-sasane}, Def. \ref{def:spacetime-hawking}) require 
		spacetime to be Lorentzian.
		As such, we will find ourselves frequently using (actually, exclusively using in practice)
		the Lorentzian metrics. We will use the Minkowski metric in special relativity, which 
		is a special subset of the Lorentzian metric, and a general Lorentzian metric for general 
		relativity.

		There are three types of vector in a vector space $V$ with a Lorentzian metric. Now
		that we have the notion of the inner product, we can use it as such:

		For any $v$ that satisfies
		\[
			g(v,v) > 0,
		\]
		then $v$ is called a \textbf{spacelike} vector. 

		Any $v$ that satisfies
		\[
			g(v,v) < 0
		\]
		is called a \textbf{timelike} vector. 

		Finally, any vector $v$ that satisfies 
		\[
			g(v,v) = 0
		\]
		is called a \textbf{lightlike} vector, or a \textbf{null vector}.

	\paragraph{Signature of the Metric}
		The summation of all the diagonal elements (the \textit{trace}, if you will),
		is called the signature of the metric. For example, for a metric
		with diagonal elements $(-1,+1,+1,+1)$, the signature is
		\[
			\sum_\mu g_{\mu\mu} = +2
		\]
		Since all the elements are $1$, we just refer to the diagonal 
		elements as their sign. Other examples of signatures are
		\[
			\sum (+,+,+,+)=4, \qquad \sum (+,-,-,-)=-2
		\]
		The signature is invariant for any basis.\footnote{This is a theorem, seach for it.}
		Sometimes in literature, the diagonals themselves are what is referred to as the signature.
		For example, many call the Lorentzian signature as $(-,+,+,+)$.

		For Lorentzian metrics, there are two conventions in the literature.
		It is presented usually (up to a trivial reording) as $(-,+,+,+)$, with signature 2. However,
		in the other convention, the Lorentzian metric is defined as a metric whose diagonal
		elements only has one $+1$, and thus its elements are $(+,-,-,-)$, and the signature is -2.
		The course adopts the convention with signature +2. 

		In the convention where they call the Lorentzian signature $(+,-,-,-)$, the definitions
		of a spacelike vector and timelike vectors are exact opposites. That is, spacelike vectors
		have $g(v,v)<0$, and timelike vectors have $g(v,v)>0$. However, there is no essential
		difference: A vector that is timelike in the $-2$ signature is also timelike in the $+2$ signature.
		This is just a difference in convention, the math remains the same. 

		The zero vector is certainly a null vector, but
		not vice versa. Many readers may only be familiar with positive definite metrics, and
		may think $v = \vb{0}$ (the zero element) whenever $g(v, v) = 0$. However, if a metric is
		Lorentzian, then $g(v, v) = 0$ does not necessarily lead to $v = 0$ (the zero element
		is unique, while there are infinitely many null vectors). Nonzero 4-dimensional null
		vectors play a significant role in relativity. For instance, it is convenient to use them
		to describe the propagation of electromagnetic waves and gravitational waves.

	\paragraph{Metric as a vector/covector map}
		Tensors are essentially multifaceted linear maps. Similar to what we discussed before,
		if you feed a tensor all its required vectors and covectors, it spits out a real number.

		A metric is a tensor of type $(0,2)$, which is a bilinear map from $V\times V\to\R$, so
		for all $u,v\in V$, we have $g(u,v)\in\R$. However, if you only give one vector, then
		the metric becomes a one-form instead
		\[
			g(v,*)\in V^*.
		\]
		
		The symbol $*$ here indicates a blank slot. This is the same if you feed the second slot, for $g(*,v)$ is also an
		element of $V^*$. In fact, we can see that they are equal for a symmetric metric.
		\[
			g(v,*) = g(v,*)
		\]
		
		That is, if you feed a metric tensor only one vector, it becomes a dual vector. When we say
		$g(v,*)$, we refer to dual.

		Given $g$, we can create $g(v,*) = g(*,v)\in V^*$ for any $v\in V$, and hence $g$ can be
		viewed as a linear map from vectors to covectors, i.e.,
		\[
			g:V\to V^*,\qquad g:\overline{v}\mapsto \underline{v}
		\]
		That is, the metric supplies the natural isomorphism between $V$ and $V^*$. A vector, $\overline{v}$, when fed 
		to a metric, becomes a dual $\underline{v}$.
		\begin{proof}
		    Exercise to the reader.
		\end{proof}
		V is naturally identified with $V^{**}$ whether or not there is a metric; if there is a metric,
		then $V$ can also be identified with $V^*$. 

		Thus, we have naturally found a way to transform a vector 
		into a covector. We can define in the notation that
		\[
			\underline{v} = g(\overline{v},*)
		\]
		This makes it clear that applying a metric on a vector alone transforms it into a covector,
	
		Returning to the manifold, we can define what is a metric tensor field is
		\begin{definition}{Metric Tensor Field}
		    A symmetric, everywhere non-degenetrate tensor field of type $(0,2)$ 
			is called a metric tensor field.
		\end{definition}		

	\paragraph{Some pointers on Abstract Index Notation}
		As we discussed last meeting, we have seen the utility of 
		abstract index notation (AIN), which we will begin to adapt for our discussions.
		For example, we can represent a tensor without specifying its type, i.e.,
		\[
			S\indices{^{ab}_c}
		\]
		Immediately shows that $S\in(2,1)$, without the specification of what the properties 
		of the tensor properties. Similarly
		\[
			v^a
		\]
		can be immediately identified as a vector (formerly we needed to specify that $v\in V$),
		and a covector can immediately also be immediately identified just with the notation.
		\[
			\omega_b
		\]

		Additionally, a tensor product is immediately identified instead of specified. In AIN, 
		a tensor product is just a juxtaposition of tensors, i.e., for a tensor $S\indices{^{ab}_c}$ and
		$T\indices{^d_e}$, the tensor product is just represented as
		\[
			S\otimes T = 
			S\indices{^{ab}_c} T\indices{^d_e} 
		\]
		If you commit to the indices which you use, i.e., the indices 
		tell you the order of where to place the place the tensors, the first rank 1 tensor
		goes to index $a$, the second to $b$, and so on, we can immediately just 
		show that:
		\[
			S\indices{^{ab}_c} 
			T\indices{^d_e} 
			=
			T\indices{^d_e} 
			S\indices{^{ab}_c} 
		\]
		It can then immediately shown, just by notation alone, the fact that the tensor product
		is noncommutative, just by showing that the order the tensors take the covectors/vectors
		are different for each tensor
		\[
			S\indices{^{ab}_c} 
			T\indices{^d_e} 
			\neq
			S\indices{^{da}_e} 
			T\indices{^b_c} 
		\]
		The advantages are that it is convenient, especially in showing 
		contractions, but the drawback is that when you deal with higher rank tensors 
		such as tensors greater than 26 dimensions, you will run out of letters.

		Feeding a tensor arguments in index-free notation. Say you have a tensor $G\in(0,2)$,
		which accepts vectors $u,v\in V$. In index free, as we have seen before, this is 
		rather cumbersome. 
		\begin{align*}
			G &= G_{\mu\nu}\dd{x^\mu}\otimes\dd{x^\nu}\\
			G(u,v) &= G_{\mu\nu}\dd{x^\mu}\otimes\dd{x^\nu}(u^\alpha \p_\alpha, v^\beta \p_\beta)\\
			G(u,v) &= G_{\mu\nu}\dd{x^\mu}(u^\alpha \p_\alpha)\dd{x^\nu}(v^\beta \p_\beta)\\
			G(u,v) &= u^\alpha v^\beta G_{\mu\nu}\delta\indices{^\mu_\alpha}\delta\indices{^\nu_\beta}\\
			G(u,v) &= u^\alpha v^\beta G_{\alpha\beta}
		\end{align*}

		In AIN, feeding a tensor just means placing a vector beside it. A contraction is then implied 
		by the matching indices.
		\[
			G(u,v) 
			=
			G_{ab} u^a v^b
		\]
		A contraction $C\indices{^i_j}T$ is identical to feeding the tensor the 
		$i$-th and $j$-th basis vector.
		\[
			C\indices{^i_j} T = T(\p_i, dx^j)
		\]
		In index-free notation, this is rather cumbersome to represent. However, in ANI, we can simplify 
		the representation. Say we have a tensor $R\indices{^h_{ijk}}$. A contraction 
		is represented as
		\begin{align*}
			R\indices{^h_{ijk}} 
			&=
			R\indices{^h_{ijk}} \delta\indices{^j_h}\\
			&=
			R\indices{^j_{ijk}} =  
			R\indices{^h_{ihk}} \\
			&=
			R\indices{_{ik}} 
		\end{align*}
		Which rather simplifies how a contraction decreases the rank of a tensor.	
		\newpage
	\subsubsection{Metric as a raising/lowering operator}
		The metric is the most importat tensor in GR and in differential geometry in general.
		It provides both a description of the curvature of the manifold (we will see more of this in the later meetings),
		and gives a sense of length and correspondence between vectors. However, another 
		utility of the metric is a map from vectors to covectors and vice versa.

		We discussed earlier that, a metric, when fed two vectors, becomes a real number.
		\[
			g_{ab} u^a v^b \in \R
		\]
		We also discussed, that a metric is symmetric by definition, i.e.,
		\[
			g_{ab} u^a u^b = g_{ba} u^a u^b
		\]

		We then discussed that a metric, when fed only one vector, becomes a covector
		\[
			g(*,\overline{u}) = \underline{u}
		\]
		It's the same kernel, but the vector has now been mapped linearly to a covector.
		It is not an isomorphisim, but a map, i.e.,
		\[
			g:V\to V^*
		\]
		or if you prefer to talk in terms of manifolds,
		\[
			g:T_p \mcal M \to T_p^* \mcal M
		\]
		The metric therefore, is a natural correspondence between the tangent space
		and cotangent space. This map is dependent on $g$ being a nondegenerate
		structure.\footnote{
			There is a similar concept in Lagrangian mechanics called the symplectic matrix, or symplectic
			structure. It is a bilinear, antisymmetric two form defined to be nondegenerate,
			which facilitates a correspondence between Lagrangian and Newtonian mechanics.
			You can look that up if you want, but combining the symplectic matrix and metric tensor forms 
			is not really common in in Physics.
		}

		In AIN, we denote that
		\[
			g(*,\overline{u}) = \underline{u} \iff u_a = g_{ab} u^b
		\]
		That is, it denotes the mapping from vectors to covectors.

		Now, because it is a one-to-one map, there must exist an inverse $g\inv$
		such that
		\[
			g\inv:V^*\to V
		\]
		or 
		\[
			g\inv:T^*_p\mcal{M}\to T_p \mcal{M}
		\]
		We can then reason that this inverse must be a $(2,0)$ tensor that
		accepts covectors and returns a vector
		\[
			(g\inv)^{ab} \omega_b = \omega^a
		\]
		\begin{proof}
		    Exercise to the reader
		\end{proof}

		By convention, we no longer write the $-1$, and just denote this
		metric inverse as
		\[
			g^{ab} \equiv (g\inv)^{ab}
		\]
		We can then show that the metric and the metric inverse when 
		applied result to an identity map, using the same kernel.
		\[
			g^{ab}g_{bc} = \Id\indices{^a_c}
		\]
		We can now then have a structure which converts from covectors to vector, i.e.,
		for a covector $u_b$, when the metric inverse is applied, becomes a vector.
		\[
			u^a = g^{ab}u_b
		\]

		The metric tensor $g_{ab}$ can then be referred to as a \textit{lowering operator}, since it lowers
		the indices of a vector, and the metric inverse $g^{ab}$ as a \textit{raising} operator, since
		it raises the indices of covectors. This can work for any rank 1 tensor in general,
		as long as the indices match upstairs and downstairs.

		\begin{align*}
			g_{ab}:\quad g_{ab}u^a = u_b &\implies \text{Lowering operator}\\
			g^{ab}:\quad g^{ab}u_a = u^b &\implies \text{Raising operator}\\
		\end{align*}

		This notation is common in GR, since raising and lowering indices is 
		performed a lot in calculations.\footnote{Raising and lowering operations are 
		also common in mathematics, but it's a lot more difficult because index-free notation
		is a pain in the ass to use.}

	\subsubsection{Symmetrization and Antisymmetrization}
		
		A tensor which accepts two vectors is in general, antisymmetric. i.e.,
		for a $(0,2)$ tensor,
		\[
			C(\overline{u},\overline{v})\in \R
		\]
		The following is generally true for 
		any arbitrary vectors $u$ and $v$.

		\[
			C(\overline{u},\overline{v})\neq C(\overline{v}, \overline{u})
		\]

		To get a symmetric result, we define an operator $\mathrm{Sym}$ such that
		\begin{equation}
			\label{eq:sym-operator}
			\mathrm{Sym}\  
			C (\overline{u},\overline{v}) 
			\equiv
			\frac{1}{2}
			\pqty{
				C (\overline{u},\overline{v}) 
				+
				C (\overline{v},\overline{u}) 
			}
		\end{equation}
		This works for $(2,0)$ tensors also
		\[
			\mathrm{Sym}\  
			C (\underline{u},\underline{v}) 
			\equiv
			\frac{1}{2}
			\pqty{
				C (\underline{u},\underline{v}) 
				+
				C (\underline{v},\underline{u}) 
			}
		\]

		In AIN, this distinction is rather easy to represent. We 
		define a symmetrization of a tensor $C_{ab}$ with a parenthesis
		on the indices. 
		\begin{equation}
			\label{eq:sym-operator-AIN}
			C_{(ab)} = \frac{1}{2}(C_{ab}+C_{ba}) 
		\end{equation}
		For rank 3 tensors, we have
		\[
			C_{(abc)} = \frac{1}{3!}
			\pqty{
				C_{abc} + 
				C_{cab} + 
				C_{bca} + 
				C_{acb} + 
				C_{bac} + 
				C_{cba}  
			}
		\]
		Where the sum of all the permulations of the indices of $C_{abc}$ are taken.
		In general, a symmetrization of a tensor is
		\[
			C_{(a1\cdots a_n)} =
			\frac{1}{n!}
			\pqty{
				C_{a_1\cdots a_n} + \cdots \text{all permutations}
			}
		\]
		The same is true for purely contravariant tensors.
		Just raise the indices.
		\[
			C^{(a1\cdots a_n)} =
			\frac{1}{n!}
			\pqty{
				C^{a_1\cdots a_n} + \cdots \text{all permutations}
			}
		\]

		If you have, for example a rank 4 covariant tensor $R_{abcd}$, and you
		only want to symmetrize between $a$ and $b$, you can group them together.
		\[
			R_{(ab)cd} 
			=
			\frac{1}{2}
			\pqty{
				R_{abcd} +
				R_{bacd}
			}
		\]
		or, if you want to symmetrize betwene $b$ and $d$, we have
		\[
			R_{a(b)c(d)} 
			=
			\frac{1}{2}
			\pqty{
				R_{abcd} +
				R_{adcb}
			}
		\]
		In general, if you want to symmetrize $m$ indices for a rank $n$
		tensor, you need to do
		\[
			R_{(a_1\cdots a_m)\cdots a_n}
			=
			\frac{1}{m!}
			\pqty{
				R_{a_1\cdots a_m} +
				\text{all permutations between } a_1\text{ to }a_m
			}
		\]
		The grouping is arbitrary, but this is the general rule.

		Symmetrization is only allowed for indices that are upstairs together, or downstairs together.
		You can only group them together with covariant or contravariant groups.
		Symmetrization between $a$ and $b$ in $C\indices{^a_b}$ is not allowed.

		If you want to symmetrize, say
		\[
			R\indices{_{(a)b}^{(c)}_d}
		\]
		you would need to perform a lowering operation for index $c$ first, then symmetrize.

		The opposite operation is called \textit{antisymmetrization}, where 
		we only want to get the antisymmetric part of a tensor. We denote it as
		\[
			C_{[ab]} = \frac{1}{2}\pqty{
				C_{ab} - C_{ba}
			} 
		\]
		We wish to permute the indices such that we only have odd permutations.
		In general
		\[
			C_{[a_1\cdots a_n]}
			=
			\frac{1}{n!}
			\sum (-1)^\pi
			C_{\pi (a)\pi (b)\pi (c)\cdots \pi (c_m)}
		\]
		where $\pi$ indicates the odd permuation.
		
		You can see how this all would be difficult in index-free notation,
		since you need to write it all out. 

		These operations are done because, as you will see in the coming weeks, we need
		to work with symmetric tensors a lot, and sometimes there are nonsymmetric tensors,
		which to manipulate mathematically, we need to symmetrize first.
		\newpage
\subsubsection{The Geodesic Equation}
	
	Metric tensors are core to the definition of the inner product.
	It defines both the inner product of a tangent space,
	\[
		g_{ab}: V\times V\to \R
	\]
	and the inner product of the cotangent space.
	\[
		g^{ab}: V^*\times V^* \to \R
	\]

	We have discussed in detail in \S\ref{subsubsec:inner-product},
	how the metric provides a sense of length of vectors.
	Now, we use that sense to provide a sense of lengths for 
	curves or lines in a manifold $\mcal M$.

	Let us begin with the intuitive notion of an arc in the Euclidean space $\R^n$.
	Suppose the parametric equation of a curve $C(t)$ in  the natural coordinate systm
	$[x,y]$ is $x= x(t)$ and $y = y(t)$, then the square length 
	of a curve segment $\dd{\ell}$ is $\dd{\ell}^2$. From Pythagoras,
	we have
	\begin{align*}
		d\ell^2 &= dx^2 + dy^2 \\
				&= \pqty{
					\pqty{
						\dv{x}{t} 
					}^2
					+
					\pqty{
						\dv{y}{t} 
					}^2
				}
				\dd{t}\\
				&=
				\pqty{
					(T^1)^2 
					+
					(T^2)^2
				}\dd{t}^2\\
				&=
				\abs{T}^2\dd{t}^2
	\end{align*}
	where $T$ is a tangent vector of $C(t)$ at $t$. From this, we get
	\[
		\dd{\ell} = \abs{T}\dd{t}
	\]
	and we define the \textit{arc length} of $C(t)$ as
	\[
		\ell = \int \abs{T}\dd{t}
	\]

	This equation can be generalized to any manifold $\mcal M$ with 
	a positive definite metric field $g$. 

	Suppose that $\gamma$ is a curve, such that
	\[
		\gamma:\R\to \mcal M
	\]
	Let the parameter of $\gamma$ be $\lambda\in I\subseteq\R$, i.e,
	the curve is $\gamma(\lambda)$.

	How do we find the vector tangent to this curve? Let us recall
	the definition of a tangent vector. Let $\dot\gamma$ be a vector such that
	\begin{equation}
		\label{eq:gamma-tanvec}
		\dot \gamma [\psi] = \dv{(\psi\circ\gamma)}{\lambda} 
	\end{equation}
	that is, we are taking the image of the curve and parametrizing it.%
	\footnote{
		A small sidenote here is that curves $\gamma(\lambda)$, for which
		$\lambda$ can take on the entire real line, i.e., $\lambda\in(-\infty,\infty)$,
		are called \textit{complete curves}. For bounded manifolds, such as torids and spheres,
		if $\lambda$ spans the entire bound, it is also called complete. Later, when we learn about singularities,
		this notion of completeness will come into play. Incompleteness means that for some curves, you cannot 
		extend the parameter to include the bounds. This happens in black holes and singularities, recall Penrose's or Hawking's incompleteness theorems.
	}
	
	In terms of a chart, $\psi\circ\gamma$, which we often just call the curve $\gamma$, we have
	\[
		\gamma = x^\alpha(\lambda).
	\]
	From which, we can define the tangent vector (from (\ref{eq:gamma-tanvec})) to be
	\begin{align*}
		\dot\gamma &= \dv{x^\alpha}{\lambda}\p_\alpha \\
				   &= \dot x^\alpha \p_\alpha
	\end{align*}

	The norm of the tangent vector $\dot \gamma$ is defined as the square root of its inner product with itself.
	\[
		\abs{\dot \gamma} = \sqrt{g(\dot\gamma,\dot\gamma)}
	\]
	Hence, the arc length of $\gamma(\lambda)$,
	from point $\lambda_1$ to $\lambda_2$, 
	can be naturally written as
	\begin{equation}
		\label{eq:gamma-arclength}
		\ell(\gamma) = \int_{\lambda_1}^{\lambda_2} \sqrt{
			g(\dot\gamma,\dot\gamma)
		}
		\dd{\lambda}
	\end{equation}
	Note here that $\ell$ is a functional. It takes a function, in this case a
	curve, evaluates it on $\lambda$, and returns a 
	real number. This is not unfamiliar, this is prominently used in Lagrangian mechanics.
	\[
		\ell : C\to \R
	\]

	Now, let's define what $g(\dot\gamma,\dot\gamma)$ is.

	Recall that $g$ is a symmetric tensor, such that
	\[
		g = g_{\mu\nu}\dd{x^\mu}\otimes\dd{x^\nu}
	\]
	We now feed this tensor the tangent vector $\dot\gamma = x^\alpha\p_\alpha$:
	\begin{align*}
		g(
			\dot{x}^\alpha\p_\alpha,
			\dot{x}^\beta\p_\beta
		)
		&=
		\dot{x}^\alpha
		\dot{x}^\beta
		g(\p_\alpha,\p_\beta)\\
		&= 
		g_{\alpha\beta} 
		\dot{x}^\alpha
		\dot{x}^\beta
	\end{align*}

	Hence, (\ref{eq:gamma-arclength}) becomes
	\[
		\ell(\gamma) = \int_{\lambda_1}^{\lambda_2} 
		\sqrt{
			g_{\alpha\beta} 
			\dot{x}^\alpha
			\dot{x}^\beta
		}
		\dd{\lambda}
	\]
	The entire expression inside the radical is just a real number, or a product of real numbers, evaluated at the curve.
	To be more explicit, we can express it as
	\begin{equation}
		\label{eq:gamma-total-arclength}
		\ell(\gamma) = \int
		\sqrt{
			g_{\alpha\beta}(x^\gamma(\lambda))
			\dv{x^\alpha}{\lambda} 
			\dv{x^\beta}{\lambda} 
		}
		\dd{\lambda}
	\end{equation}
	Note here that we have expressed the metric components $g_{\alpha\beta}$ as a function.
	It is a function that depends on the coordinate representation of the curve $\gamma=x^\gamma(\lambda)$,
	which is also dependent on the parameter $\lambda$

	Now, if we aim to determine the total arc length of the curve, then we are done. Eq. (\ref{eq:gamma-total-arclength})
	would suit your purposes just fine.

	However, what if we want more? What if we want to find the extremizations
	of the arclength? The maximal or minimal paths? Then we perform the tools used in variational principles, i.e.,
	the principle of least action.

	Recall Lagrangian mechanics. Given a functional $S$, such that:
	\[
		S[q] = \int_{t_1}^{t_2} L(q,\dot{q},t)\dd{t}
	\]
	where $L$ is the Lagrangian,
	the physical trajectory $q(t)$ is the one for which the action is \textit{stationary}
	under small variations $\delta q$ that vanish at the end points, i.e.,
	\[
		\var{S} = 0
	\]

	When you carry out the variation $\var{S}$ and set it to zero, we get the Euler-Lagrange equations\footnote{
		This is really the Euler equation. Euler formulated this, but it is Lagrange that gave it physical meaning.
		However, the concept is really only by Euler. We only call it Euler-Lagrange in the context of mechanics.
	}
	\begin{equation}
		\label{eq:euler-lagrange}
	    \dv{}{\lambda}
		\pqty{
			\pdv{L}{\dot q^\mu}  
		}
		=
		\pdv{L}{q^\mu}
	\end{equation}

	To make our lives easier, consider (\ref{eq:gamma-total-arclength}). Note that the square root, $\sqrt{x}$,
	is a monotonically increasing function.  As such, if we do minimization, it really has no 
	effect on the stationary path, and minimizing the length is equivalent to minimizing the argument of the radical, i.e.,
	minimizing $\sqrt{x}$ has the same effect as minimizing $x$
	Therefore, we might as well just use $x$.
	\begin{equation}
		\label{eq:gamma-action}
		S
		= \int
		{
			g_{\alpha\beta}(x^\gamma(\lambda))
			\dv{x^\alpha}{\lambda} 
			\dv{x^\beta}{\lambda} 
		}
		\dd{\lambda}
	\end{equation}
	So we have this ``action'' like term, $S$, and a ``Lagrangian'' like term, 
	$
			L=g_{\alpha\beta}(x^\gamma(\lambda))
			\dv{x^\alpha}{\lambda} 
			\dv{x^\beta}{\lambda} 
	$. This can be easily extremized using the Euler-Lagrange formulation.
	We now perform variation on $S$, requiring $\var{S}= 0$.

	We solve for each of the terms of (\ref{eq:euler-lagrange})

	For the RHS, we have
	\begin{align*}
	    \pdv{L}{x^\mu} 
		&= 
		\pdv{g_{\alpha\beta}}{\lambda^\mu}
		\dot{x}^\alpha
		\dot{x}^\beta\\
		&=
		\p_\mu g_{\alpha\beta}
		\dot{x}^\alpha
		\dot{x}^\beta
	\end{align*}

	The term $\displaystyle \pdv{g_{\alpha\beta}}{x^\mu}$ comes up
	so often in GR and Differential geometry that we make a shorthand notation for it.\footnote{
		Sir uses $g_{\alpha\beta,\mu}$, and he did for the lecture.
		However, I prefer to use $\p_\mu g_{\alpha\beta}$, since it is more
		indicative of the derivative operation. That is what we will be using 
		for the rest of this notes. Just bear in mind it is equivalent.
	}
	\[
		\pdv{g_{\alpha\beta}}{x^\mu} \equiv \p_\mu g_{\alpha\beta} \equiv g_{\alpha\beta,\mu} 
	\]

	For the LHS, we have
	\begin{align*}
		\pdv{L}{\dot{x}^\mu}
		&= 
		g_{\alpha\beta} \pdv{}{\dot{x}^\mu}  
		\pqty{
				\dot{x}^\alpha \dot{x^\beta}
			}\\
		&= 
		g_{\alpha\beta} 
		\pqty{
			\dot{x^\beta}
			\pdv{
				\dot{x}^\alpha 
			}{
				\dot{x}^\mu
			} 
			+
			\dot{x^\alpha}
			\pdv{
				\dot{x}^\beta 
			}{
				\dot{x}^\mu
			} 
		}\\
		&= 
		g_{\alpha\beta} 
		\pqty{
			\dot{x^\beta}
			\delta\indices{^\alpha_\mu}
			+
			\dot{x^\alpha}
			\delta\indices{^\beta_\mu}
		}\\
		&= 
		g_{\alpha\beta} 
		\dot{x^\beta}
		\delta\indices{^\alpha_\mu}
		+
		g_{\alpha\beta} 
		\dot{x^\alpha}
		\delta\indices{^\beta_\mu}\\
		&= 
		g_{\mu\beta} 
		\dot{x^\beta}
		+
		g_{\alpha\mu} 
		\dot{x^\alpha}
	\end{align*}
	Since the metric is symmetrical, i.e., $g_{\mu\beta}= g_{\beta\mu}$, we can
	use the symmetry argument to just show
	\[
		g_{\mu\beta}\dot{x}^\beta = g_{\beta\mu}\dot{x}^\beta
	\]
	And we rename the second dummy index to $\alpha$, hence
	\[
		g_{\mu\beta}\dot{x}^\beta = g_{\mu\alpha}\dot{x}^\alpha
	\]
	So the two terms are identical, and their sum is just twice of one of them. We choose the $\beta$ term.
	\[
		\pdv{L}{\dot{x}^\mu}
		= 
		2g_{\mu\beta}\dot{x}^\beta
	\]
	We then do the derivation operation
	\begin{align*}
	    \dv{}{\lambda}
		\pqty{
			\pdv{L}{\dot{x}^\mu}
		}
		&= 
		2\dv{g_\mu\beta}{\lambda}\dot{x}^\beta + 2g_{\mu\beta}\dv{\dot{x}^\beta}{\lambda}\\
		&= 
		2\pdv{g_{\mu\beta}}{x^\alpha}\dv{x^\alpha}{\lambda}\dot{x}^\beta + 2g_{\mu\beta}\dv{\dot{x}^\beta}{\lambda}\\
		&= 
		2\p_\alpha{g_{\mu\beta}}\dot{x}^\alpha \dot{x}^\beta + 2g_{\mu\beta}\ddot{x}^\beta
	\end{align*}
	Thus, equating the RHS and LHS together, and dividing by 2 on both sides, we have
	\[
		\frac{1}{2} 
		\p_\mu g_{\alpha\beta}
		\dot{x}^\alpha
		\dot{x}^\beta
		=
		\p_\alpha{g_{\mu\beta}}\dot{x}^\alpha \dot{x}^\beta 
		+ 
		g_{\mu\beta}\ddot{x}^\beta
	\]
	This is just a plain, garden variety, system of linear ordinary differential equations.
	We rewrite it to isolate the second order term
	\[
		g_{\mu\beta}\ddot{x}^\beta
		=
		\frac{1}{2} 
		\p_\mu g_{\alpha\beta}
		\dot{x}^\alpha
		\dot{x}^\beta
		-
		\p_\alpha{g_{\mu\beta}}\dot{x}^\alpha \dot{x}^\beta 
	\]
	We rewrite this in a rather strange form, but you will see its utility later
	\[
		g_{\mu\beta}\ddot{x}^\beta
		=
		-
		\frac{1}{2} 
		\pqty{
			2
			\p_\alpha{g_{\mu\beta}}
			\dot{x}^\alpha 
			\dot{x}^\beta 
			-
			\p_\mu g_{\alpha\beta}
			\dot{x}^\alpha
			\dot{x}^\beta
		}
	\]
	We can pull out the common terms
	\[
		g_{\mu\beta}\ddot{x}^\beta
		=
		-
		\frac{1}{2} 
		\pqty{
			2
			\p_\alpha{g_{\mu\beta}}
			-
			\p_\mu g_{\alpha\beta}
		}
		\dot{x}^\alpha 
		\dot{x}^\beta 
	\]
	Really, we're done here. This is the equation we are looking for, and in fact it is easier to work with.
	However, we can make the equation look neater.

	Here, we again exploit the symmetries of the system. Since
	this is just ordinary multiplication, it is symmetric.
	Because $\alpha$ and $\beta$ are dummy indices, we can interchange them.
	That is,
	\[
		\p_\alpha
		g_{\mu\beta} 
		\dot{x}^\alpha 
		\dot{x}^\beta 
		=
		\p_\beta
		g_{\mu\alpha}
		\dot{x}^\beta 
		\dot{x}^\alpha 
	\]
	And because ordinary multiplication commutes,
	\[
		\dot{x}^\alpha 
		\dot{x}^\beta 
		=
		\dot{x}^\beta 
		\dot{x}^\alpha 
	\]
	Therefore,
	\[
		\p_\alpha
		g_{\mu\beta} 
		\dot{x}^\alpha 
		\dot{x}^\beta 
		=
		\p_\beta
		g_{\mu\alpha}
		\dot{x}^\alpha 
		\dot{x}^\beta 
	\]
	We can average these two equal expressions,
	\[
		\p_\alpha
		g_{\mu\beta} 
		\dot{x}^\alpha 
		\dot{x}^\beta 
		=
		\frac{1}{2}
		\pqty{
			\p_\alpha
			g_{\mu\beta} 
			\dot{x}^\alpha 
			\dot{x}^\beta 
			+
			\p_\beta
			g_{\mu\alpha}
			\dot{x}^\alpha 
			\dot{x}^\beta 
		}
		\dot{x}^\alpha 
		\dot{x}^\beta 
	\]
	Thus we can separate the $2 \p_\alpha{g_{\mu\beta}}$ term to have
	\[
		g_{\mu\beta}\ddot{x}^\beta
		=
		-
		\frac{1}{2} 
		\pqty{
			\p_\alpha{g_{\mu\beta}}
			+
			\p_\beta{g_{\mu\alpha}}
			-
			\p_\mu g_{\alpha\beta}
		}
		\dot{x}^\alpha 
		\dot{x}^\beta 
	\]
	Now, we want to isolate the $\ddot{x}^\beta$ term to be on its own.
	Notice that when we apply the inverse metric,
	\[
		g^{\sigma\mu}
		g_{\mu\beta}
		=
		\delta\indices{^\sigma_\beta}
	\]
	Hence
	\begin{align*}
		g^{\sigma\mu}
		g_{\mu\beta}
		\ddot{x}^\beta
		&=
		-
		\frac{1}{2} 
		g^{\sigma\mu}
		\pqty{
			\p_\alpha{g_{\mu\beta}}
			+
			\p_\beta{g_{\mu\alpha}}
			-
			\p_\mu g_{\alpha\beta}
		}
		\dot{x}^\alpha 
		\dot{x}^\beta \\
		\delta\indices{^\sigma_\beta}
		\ddot{x}^\beta
		&=
		-
		\frac{1}{2} 
		g^{\sigma\mu}
		\pqty{
			\p_\alpha{g_{\mu\beta}}
			+
			\p_\beta{g_{\mu\alpha}}
			-
			\p_\mu g_{\alpha\beta}
		}
		\dot{x}^\alpha 
		\dot{x}^\beta \\
		\ddot{x}^\sigma
		&=
		-
		\frac{1}{2} 
		g^{\sigma\mu}
		\pqty{
			\p_\alpha{g_{\mu\beta}}
			+
			\p_\beta{g_{\mu\alpha}}
			-
			\p_\mu g_{\alpha\beta}
		}
		\dot{x}^\alpha 
		\dot{x}^\beta 
	\end{align*}

	This is the form we are looking for. To tidy up, the term in front of 
	$
		\dot{x}^\alpha 
		\dot{x}^\beta 
	$ is so
	common in GR and differential geometry, that it has a name. It is called the Christoffel symbols of the first kind.
	\begin{definition}{Christoffel Symbols}
		For a manifold $\mcal M$ with metric $g_{\mu\nu}$,
		the Christoffel symbols $\Gamma^\sigma_{\alpha\beta}$ are 
		the connection coefficients which encode how coordinate 
		basis vectors vary as one moves along the manifold.
		\begin{equation}
			\label{eq:christoffel-symbol}
			\Gamma^{\sigma}_{\alpha\beta}
			= 
			\frac{1}{2} 
			g^{\sigma\mu}
			\pqty{
				\p_\alpha{g_{\mu\beta}}
				+
				\p_\beta{g_{\mu\alpha}}
				-
				\p_\mu g_{\alpha\beta}
			}
		\end{equation}
	\end{definition}
	
	Hence, we can now rewrite our final result with this symbol, and 
	we derive the the extremal path
	in a manifold, the famous geodesic equation.

	\begin{theorem}{Geodesic equation}
		A free particle follows
		the straightest possible path in a manifold $\mcal{M}$ 
		(a geodesic) $x^\sigma(\lambda)$ satisfying

		\begin{equation}
			\label{eq:geodesic-eqn}
			\ddot{x}^\sigma
			+
			\Gamma^{\sigma}_{\alpha\beta}
			\dot{x}^\alpha 
			\dot{x}^\beta 
			=
			0
		\end{equation}
	\end{theorem}

	If we expand it to explicitly show the dependencies on the 
	curve, we have
	\begin{equation}
		\label{eq:geodesic-eqn-explicit}
		\dv[2]{x^\sigma}{\lambda} 
		+ 
		\Gamma^\sigma_{\alpha\beta}(x^\sigma(\lambda)) 
		\dv{x^\alpha}{\lambda} 
		\dv{x^\beta}{\lambda} 
		=
		0
	\end{equation}
	To be specific, it is the extremal path on Riemannian 
	and semi-Riemannian manifolds.

	This does not immediately imply the shortest or minimal route
	in a manifold. To do that, you have to check the second derivative.

	In Riemannian mafolds, where all the elements of the metric is positive,
	it is indeed the minimal path. 
	However, for Lorentzian manifolds, it is the maximal path, the path 
	where proper time will be maximized.

	\paragraph{Significance}

	Solving geodesics have a lot of implications. It 
	has a lot of physical significances. Firstly, 
	for normal geometry or for positive definite manifolds, 
	it is the shortest or straightest 
	possible path between two points, the most optimized
	curve between them. Second, it is the curve where
	vectors remain parallel to their starting point (this is known
	as a parallel transport).
	However, most significantly, it is the special curve
	that the world line of freely falling bodies follow.

	To get the path of inertial observers, solve the geodesic equation.

	To get the orbit of planets, such as the orbits of Earth, or Mercury.
	use the Schwarzschild metric and solve the geodesic equation.
	All inertial observers follow the path of maximuzed proper time,
	and to obtain that path requirees solving the geodesic equation.

	Solving the equation is very hard. Firstly, it is obviously
	nonlinear, so you have elements of chaos theory affecting your results,
	so small differences in initial conditions lead to large differences
	in time propagation.
	Geodesics in general space time, i.e., arbitrary metrics, are 
	nigh impossible to solve. Most of the times, it is solved numerically.

	As such, we will learn techniques in the following weeks,
	that exploit conserved quantities and symmetries in the metric
	that allow us to solve for it generally. We will learn Killing vectors,
	which exploit the symmetries in space time and give conserved
	quantities, isometries, such as symmetries of the metric.
	
	For a general metric with no symmetries, solving the 
	geodesic equation is impossible analytically. That's 
	why people resort to numerical relativity\footnote{The 
	objectively correct path} to do GR, because analytic solutions
	are rare and those that are trivial (or at least relatively easy),
	such as the Minkowski or Schwarzschild metrics have all been solved
	analytically. 

	Normally today, especially in higher level research, all the work we do 
	is done using \textsc{Mathematica}, or some other 
	computer program, to solve for the geodesic equation,
	for the metric, and for the Christoffel symbols, for it is simply too 
	cumbersome to solve for them by hand analytically.
	
	\begin{remark}
	    The Christoffel symbol is not a tensor!
		
		\hspace{1em}
		While it is indexed, i.e.,
		\[
			\Gamma^\sigma_{\alpha\beta}
		\]
		However, it does not transform like a tensor. 
		The components of a Christoffel tensor is 
		known as a connection.

		\hspace{1em}Here we learn that not all indexed quantities are 
		tensors. For example, the Levi-Civita symbol 
		$\varepsilon\indices{_{ij}^k}$
		is not a tensor. It is a pseudotensor.
		This is a similar case with the Christoffel symbols.
	\end{remark}

	\begin{remark}
		There is a dual version of the geodesic equation. The equation describes a vector acceleration. To convert, all you need to do is
		multiply a lowering operator for each vector, and you will have the covector acceleration.
		\begin{equation}
			\label{eq:geodesic-eqn-dual}
			\ddot{x}_\sigma 
			+
			\Gamma^{\alpha\beta}_\sigma
			\dot{x}_\alpha
			\dot{x}_\beta
			=0
		\end{equation}
	\end{remark}
	There is a very similar class of curves to geodesic curves, called autoparallel curves, which 
	are curves that parallel transport tangent vectors. They are governed
	by the autoparallel equation, which we will introduced later. It is very similar to the geodesic equation. They
	are natural connections on manifolds, which require the coefficients
	of the connection to be Christoffel. 

	We are motivating this because we are setting up the path to describe what a parallel transport of 
	a tangent vector is, which is going to be important later on.
	\begin{question}{If a metric is not dependent on the parameter $\lambda$, does it mean it is 0 all throughout?}

		No, not quite. 
		Consider the Scwharschild metric,
		\[
			\dd{s}^2 = 
			-\pqty{
				1 - \frac{2GM}{rc^2} 
			}\dd{t}^2
			+
			\pqty{
				1 - 
				\frac{2GM}{rc^2} 
			}\inv \dd{r}^2
			+ 
			r^2 \pqty{
				\dd{\theta}^2 + 
				\sin^2\theta\dd{\phi}
			}
		\]
		The chart is the Schwarzschild coordinates $(t,r,\theta,\phi)$
		Notice that for the $\dd{t}$ term, there is no time dependence.
		However, it does not mean that the result is zero.
		\[
			\pqty{
				\pdv{L}{\dot t}\neq 0  
			}  
		\]
		Rather, it implies that there is a constant.
		\[
			{
				\pdv{L}{\dot t} = \text{constant}
			}  
		\]
		In this case, you will find that for this term, the constant
		is equal to $E/m$, the mass per energy unit of the particle.
		Similarly, for $\pdv{L}{\phi}$, which has no $\phi$ dependence,
		you will find that 
		\[
			\pdv{L}{\dot \phi} = \text{constant} 
		\]
		In this case it is the angular momentum per unit mass.
		These constants describe the conserved quantities for a metric.
	\end{question}
